% !TEX root = ../main.tex
\chapter{中心极限定理}
\label{ch:clt}

\noindent\emph{用到的前置工具：随机变量、分布与期望（概率与测度部分）；独立性与多元正态（『多维分布、独立性与变量变换』一章）；依分布收敛（见前文『收敛模式』一章）；大数定律（见前文『大数定律』一章）；Taylor 展开（多元微积分部分）。}
\medskip

上一章的大数定律告诉我们样本均值会\emph{收敛到}总体均值——它定住了估计量的``靶心''，却没说估计量在靶心周围如何抖动。可计量经济学真正赖以生存的，恰恰是后者：一个置信区间要多宽、一个 $t$ 统计量该和什么分布比、一个 Wald 检验的临界值从哪来，全都取决于估计量\emph{围绕真值的波动形态}。中心极限定理（central limit theorem，CLT）正是回答这个问题的工具。它给出一个出人意料地干净的答案：只要把样本均值减去其期望、再按 $\rtn$ 放大，无论原始数据服从什么分布，这个标准化后的量都会收敛到\emph{同一个}极限——标准正态。

这件事的份量怎么强调都不过分。大数定律说``$\bar X_n$ 趋于 $\mu$'',是一句关于``位置''的话；中心极限定理说``$\rtn(\bar X_n-\mu)$ 趋于 $\cN(0,\sigma^2)$'',是一句关于``形状''的话——而正是这个形状，让我们能给估计量配上标准误、画出置信区间、算出 $p$ 值。本章先说清这件工具解决什么问题、$\rtn$ 这个标度从何而来，再用特征函数法把``为什么极限恰是正态''讲透，然后陈述放松了同分布假设的版本（Lyapunov / Lindeberg）与多元版本，最后指明它在计量推断里的几处主战场。

\section{问题：估计量围绕真值如何波动}
\label{sec:clt-motivation}

设 $X_1,\dots,X_n$ 独立同分布（iid），均值 $\E{X_i}=\mu$、方差 $\Var{X_i}=\sigma^2\in(0,\infty)$。样本均值
\[
    \bar X_n=\frac1n\sumi X_i
\]
是 $\mu$ 的最自然估计。大数定律给出 $\bar X_n\pto\mu$：随 $n\to\infty$，整条抽样分布会塌缩到 $\mu$ 这一个点。但``塌缩''本身没法用来做推断——若极限是一个退化于 $\mu$ 的点质量，我们就无从谈论``$\bar X_n$ 偏离 $\mu$ 多远算正常''。

要看清波动的形状，得把这个不断收缩的分布\emph{放大}到一个非退化的尺度上。先算一下该放大多少。由独立性，
\[
    \Var{\bar X_n}=\frac{1}{n^2}\sumi\Var{X_i}=\frac{\sigma^2}{n}.
\]
方差以 $1/n$ 的速率趋于 $0$，故标准差以 $1/\rtn$ 的速率趋于 $0$。要让放大后的量有一个固定（不为 $0$ 也不发散）的方差，就得乘以 $\rtn$：
\[
    \Var{\rtn\,(\bar X_n-\mu)}
    = n\cdot\Var{\bar X_n}
    = n\cdot\frac{\sigma^2}{n}
    = \sigma^2 .
\]
这正是 $\rtn$ 标度的来历——它不是凑出来的魔法常数，而是\emph{恰好抵消方差收缩速率}的那个唯一标度。乘少了（比如只乘 $n^{1/4}$）量仍塌缩到 $0$，乘多了（比如乘 $n$）量发散到无穷；唯有 $\rtn$ 把分布定格在一个固定的、可供推断的非退化形状上。

\state{$\rtn$ 标度是 CLT 的中枢}{
样本均值的方差是 $\sigma^2/n$。为把不断收缩的抽样分布稳定到一个非退化极限，标准化\emph{必须}按 $\rtn$ 放大：$\dfrac{\bar X_n-\mu}{\sigma/\rtn}=\rtn\,\dfrac{\bar X_n-\mu}{\sigma}$。计量里一切以 $\rtn(\hat\theta-\theta_0)$ 形式出现的渐近结果，量纲上的根据都在这里：估计量以 $\Op(n^{-1/2})$ 的速率收敛，故须用 $\rtn$ 校准才能看清其极限分布。
}

剩下的问题只有一个：放大后这个量\emph{长什么形状}？中心极限定理的回答是——总是同一个钟形。

\section{Lindeberg–Lévy 中心极限定理}
\label{sec:clt-iid}

\thm{Lindeberg–Lévy 中心极限定理（iid 情形）}{
设 $X_1,X_2,\dots$ 独立同分布，$\E{X_i}=\mu$，$\Var{X_i}=\sigma^2\in(0,\infty)$。则
\[
    \rtn\,(\bar X_n-\mu)\ \dto\ \cN(0,\sigma^2),
    \qquad\text{等价地}\qquad
    \frac{\rtn\,(\bar X_n-\mu)}{\sigma}\ \dto\ \cN(0,1).
\]
}

这里只要求``有限的二阶矩''——存在均值与方差，别无其它。不需要数据本身正态，不需要对称，不需要连续；离散的、偏斜的、重尾但二阶矩仍有限的分布，统统适用。结论里那个正态，是从这些五花八门的原始分布里\emph{涌现}出来的，而非假设进去的。

图~\ref{fig:clt-1} 把这件事画了出来：左边是一个明显右偏、毫不像正态的原始总体；右边是其样本均值标准化后的抽样分布，随 $n$ 增大一步步贴近黑色的 $\cN(0,1)$ 钟形。$n=2$ 时仍带着原分布的偏斜，$n=5$ 时偏斜大幅减弱，到 $n=30$ 已很接近极限钟形（仔细看仍留一点右偏的尾巴——指数总体收敛偏慢）。

\begin{figure}[h]
\centering
\begin{tikzpicture}[>=Stealth]
  % ============ 左面板：原始（偏斜）总体分布 ============
  \begin{scope}
    \draw[->] (-0.2,0) -- (4.6,0) node[right] {$x$};
    \draw[->] (0,-0.2) -- (0,2.4) node[above] {密度};
    \draw[thick,orange!80!black] plot coordinates
      {(0.000,1.800) (0.165,1.549) (0.330,1.333) (0.495,1.148) (0.660,0.988)
       (0.825,0.850) (0.990,0.732) (1.155,0.630) (1.320,0.542) (1.485,0.467)
       (1.650,0.402) (1.815,0.346) (1.980,0.298) (2.145,0.256) (2.310,0.220)
       (2.475,0.190) (2.640,0.163) (2.805,0.141) (2.970,0.121) (3.135,0.104)
       (3.300,0.090) (3.465,0.077) (3.630,0.066) (3.795,0.057) (3.960,0.049)
       (4.125,0.042) (4.290,0.036)};
    \node[orange!75!black,anchor=west] at (1.7,1.55) {原始总体 $X$};
    \node[anchor=north] at (2.35,-0.35) {\small（右偏，远非正态）};
  \end{scope}
  % ============ 右面板：标准化抽样分布逐步逼近正态 ============
  \begin{scope}[xshift=9.6cm]
    \draw[->] (-3.4,0) -- (3.5,0) node[right] {$z$};
    \draw[->] (0,-0.2) -- (0,3.5) node[above] {密度};
    \foreach \x/\lab in {-2/$-2$, 2/$2$} \draw (\x,0.05)--(\x,-0.05) node[below]{\lab};
    \draw[very thick,black] plot coordinates
      {(-3.00,0.027) (-2.80,0.047) (-2.60,0.081) (-2.40,0.134) (-2.20,0.213)
       (-2.00,0.324) (-1.80,0.474) (-1.60,0.666) (-1.40,0.898) (-1.20,1.165)
       (-1.00,1.452) (-0.80,1.738) (-0.60,1.999) (-0.40,2.210) (-0.20,2.346)
       (0.00,2.394) (0.20,2.346) (0.40,2.210) (0.60,1.999) (0.80,1.738)
       (1.00,1.452) (1.20,1.165) (1.40,0.898) (1.60,0.666) (1.80,0.474)
       (2.00,0.324) (2.20,0.213) (2.40,0.134) (2.60,0.081) (2.80,0.047)
       (3.00,0.027)};
    \draw[thick,red!70!black,densely dotted] plot coordinates
      {(-1.40,0.167) (-1.30,1.166) (-1.20,1.899) (-1.10,2.418) (-1.00,2.767)
       (-0.90,2.982) (-0.80,3.092) (-0.70,3.121) (-0.60,3.089) (-0.50,3.011)
       (-0.40,2.900) (-0.30,2.766) (-0.20,2.617) (-0.10,2.459) (0.00,2.297)
       (0.10,2.135) (0.20,1.976) (0.30,1.821) (0.40,1.673) (0.50,1.533)
       (0.60,1.400) (0.70,1.276) (0.80,1.160) (0.90,1.053) (1.00,0.953)
       (1.10,0.862) (1.20,0.778) (1.30,0.701) (1.40,0.631) (1.50,0.567)
       (1.60,0.509) (1.80,0.409) (2.00,0.328) (2.20,0.261) (2.40,0.208)
       (2.60,0.165) (2.80,0.130) (3.00,0.103)};
    \draw[thick,blue!65!black,dashed] plot coordinates
      {(-2.00,0.026) (-1.80,0.191) (-1.60,0.552) (-1.40,1.053) (-1.20,1.588)
       (-1.00,2.057) (-0.80,2.396) (-0.60,2.581) (-0.40,2.618) (-0.20,2.531)
       (0.00,2.354) (0.20,2.120) (0.40,1.859) (0.60,1.593) (0.80,1.337)
       (1.00,1.104) (1.20,0.897) (1.40,0.719) (1.60,0.570) (1.80,0.446)
       (2.00,0.346) (2.20,0.266) (2.40,0.203) (2.60,0.154) (2.80,0.116)
       (3.00,0.086)};
    \draw[thick,green!50!black] plot coordinates
      {(-3.00,0.003) (-2.80,0.011) (-2.60,0.028) (-2.40,0.067) (-2.20,0.139)
       (-2.00,0.259) (-1.80,0.438) (-1.60,0.680) (-1.40,0.978) (-1.20,1.312)
       (-1.00,1.650) (-0.80,1.960) (-0.60,2.207) (-0.40,2.367) (-0.20,2.427)
       (0.00,2.387) (0.20,2.258) (0.40,2.061) (0.60,1.819) (0.80,1.556)
       (1.00,1.292) (1.20,1.043) (1.40,0.821) (1.60,0.630) (1.80,0.473)
       (2.00,0.347) (2.20,0.250) (2.40,0.176) (2.60,0.122) (2.80,0.083)
       (3.00,0.055)};
    \node[anchor=west,red!70!black] at (1.15,2.75) {\small$n=2$};
    \node[anchor=west,blue!65!black] at (1.15,2.45) {\small$n=5$};
    \node[anchor=west,green!50!black] at (1.15,2.15) {\small$n=30$};
    \node[anchor=west,black] at (1.15,1.85) {\small$\cN(0,1)$};
  \end{scope}
\end{tikzpicture}
\caption{即便原始总体明显偏斜（左），其样本均值的标准化抽样分布也随 $n$ 增大逐步逼近标准正态钟形（右）。这里原始总体取指数分布；$n=30$ 的曲线已很接近 $\cN(0,1)$（仍带一丝可见的右偏）。}
\label{fig:clt-1}
\end{figure}

\rmkb[CLT 与 LLN 是同一现象的两个尺度]{
大数定律与中心极限定理说的是\emph{同一个}收缩过程，只是看的放大倍率不同。不放大（看 $\bar X_n$ 本身）：分布塌缩到点 $\mu$，这是 LLN。放大 $\rtn$ 倍（看 $\rtn(\bar X_n-\mu)$）：塌缩被恰好抵消，露出一个固定的钟形，这是 CLT。所以 CLT 常被称作 LLN 的``二阶''细化：LLN 给收敛的\emph{位置}，CLT 给收敛的\emph{速率与涨落形状}。
}

\section{为什么极限恰是正态：特征函数法}
\label{sec:clt-cf}

要证明``标准化和趋于正态'',最锋利的工具是特征函数。它把``独立随机变量相加''这件在密度层面要做卷积的麻烦事，变成``特征函数相乘''这件代数上的易事；而依分布收敛，又恰好可由特征函数的逐点收敛来判别（Lévy 连续性定理）。两者一合，证明几乎是自动展开的。

\defn{特征函数}{
随机变量 $X$ 的\textbf{特征函数}（characteristic function）为
\[
    \vp_X(t)=\E{e^{\,itX}}=\E{\cos tX}+i\,\E{\sin tX},
    \qquad t\in\RR .
\]
它对一切随机变量都有定义（被积函数有界，$\abs{e^{itX}}=1$），且唯一地决定分布。
}

特征函数有两条性质是这整套证明的全部支柱。

\fact{特征函数的两条关键性质}{
设 $X,Y$ 独立，$a,b$ 为常数。则
\begin{enumerate}
\item \textbf{独立和化为乘积：}$\ \vp_{X+Y}(t)=\vp_X(t)\,\vp_Y(t)$；更一般地，对线性变换 $\vp_{aX+b}(t)=e^{\,itb}\,\vp_X(at)$。
\item \textbf{矩生成 Taylor 系数：}若 $\E{\abs{X}^2}<\infty$，则在 $t=0$ 附近
\[
    \vp_X(t)=1+it\,\E{X}-\tfrac12\,t^2\,\E{X^2}+o(t^2).
\]
\end{enumerate}
}

第一条来自独立时期望可分解：$\E{e^{it(X+Y)}}=\E{e^{itX}e^{itY}}=\E{e^{itX}}\E{e^{itY}}$。第二条来自把 $e^{itX}$ 在 $t=0$ 处对 $t$ 做 Taylor 展开 $e^{itX}=1+itX-\tfrac12 t^2X^2+\cdots$ 再取期望（有限二阶矩保证逐项取期望与余项控制成立）。把这两条性质对准``标准化的 iid 和'',正态就会自己长出来。

\thm{Lévy 连续性定理}{
设 $\set{Z_n}$ 为随机变量列，$\vp_{Z_n}$ 为其特征函数。若存在函数 $\vp$，它在 $t=0$ 处连续，且对每个 $t$ 都有 $\vp_{Z_n}(t)\to\vp(t)$，则 $\vp$ 是某随机变量 $Z$ 的特征函数，且 $Z_n\dto Z$。
}

这条定理是把``特征函数收敛''翻译成``分布收敛''的桥。它的深刻处（即``逐点收敛的极限确实是一个特征函数''）属于纯分析，本书按使用者的需要把它当作已知工具，重点放在如何用它收尾。有了它，证明 $Z_n\dto\cN(0,1)$ 就只剩一件事：证明 $\vp_{Z_n}(t)\to e^{-t^2/2}$——因为 $e^{-t^2/2}$ 恰是标准正态的特征函数。

\thmp{Lindeberg–Lévy CLT（特征函数法）}{
不妨设 $\mu=0$、$\sigma^2=1$（否则先做 $X_i\mapsto(X_i-\mu)/\sigma$，这不改变结论的形式）。要证 $Z_n=\dfrac{1}{\rtn}\sumi X_i\dto\cN(0,1)$。
}{
记 $\vp(t)=\vp_{X_i}(t)$ 为单个 $X_i$ 的特征函数（iid，故各项相同）。由\textbf{性质 1}，$Z_n=\frac{1}{\rtn}\sum_i X_i$ 是独立项之和再缩放，其特征函数为
\[
    \vp_{Z_n}(t)
    =\prodi \vp_{X_i}\!\pr{\frac{t}{\rtn}}
    =\br{\vp\!\pr{\frac{t}{\rtn}}}^{n}.
\]
由\textbf{性质 2}，并代入 $\E{X_i}=0$、$\E{X_i^2}=\sigma^2=1$，把自变量取作 $t/\rtn$（当 $n\to\infty$ 时它趋于 $0$，故 Taylor 展开适用）：
\[
    \vp\!\pr{\frac{t}{\rtn}}
    =1+i\,\frac{t}{\rtn}\cdot 0
    -\frac12\,\frac{t^2}{n}\cdot 1
    +o\!\pr{\frac1n}
    =1-\frac{t^2}{2n}+o\!\pr{\frac1n}.
\]
于是
\[
    \vp_{Z_n}(t)=\br{\,1-\frac{t^2}{2n}+o\!\pr{\frac1n}\,}^{n}.
\]
对固定的 $t$，用极限 $\bigl(1+\tfrac{a_n}{n}\bigr)^n\to e^{a}$（当 $a_n\to a$）——这里 $a_n=-\tfrac{t^2}{2}+o(1)\to-\tfrac{t^2}{2}$——得
\[
    \vp_{Z_n}(t)\ \longrightarrow\ e^{-t^2/2}\qquad(\forall\,t).
\]
极限 $e^{-t^2/2}$ 在 $t=0$ 处连续，正是 $\cN(0,1)$ 的特征函数。由 Lévy 连续性定理，$Z_n\dto\cN(0,1)$。
}

这条证明把``为什么是正态''说得一清二楚。看那个关键展开：单项特征函数的一阶项 $it\E{X}$ 因中心化而消失，留下的最低阶信息是\emph{二阶项} $-\tfrac12 t^2\E{X^2}$——也就是说，标准化之后，原始分布的所有细节（偏度、峰度、是离散还是连续）都退到 $o(t^2)$ 的高阶项里，唯有``方差''这一个数留在主导项上。$n$ 个这样的因子相乘（特征函数取 $n$ 次幂），把 $1-\tfrac{t^2}{2n}$ 复利式地累积成 $e^{-t^2/2}$。

\state{正态是``唯一记得均值与方差、忘掉其余''的吸引子}{
标准化抹平了一阶项，于是不同分布在 $\rtn$ 标度下的差异只剩二阶（方差）；而 $\br{1-\tfrac{t^2}{2n}}^n\to e^{-t^2/2}$ 把这唯一残存的二阶信息卷成了高斯。正因为``只有方差幸存'',形形色色的原始分布才统统被吸进同一个极限——正态分布是独立小扰动累加的\emph{吸引子}。这也解释了经验世界里钟形为何无处不在：任何由大量独立小因素叠加而成的量（测量误差、身高、资产收益的短期累积）都倾向于近似正态。
}

图~\ref{fig:clt-2} 把这个``二阶吻接 $\Rightarrow$ 高斯''的机理画了出来：任意有限二阶矩分布的标准化特征函数 $\vp_{Z_n}(t)$，在 $t=0$ 处都与同一条抛物线 $1-\tfrac{t^2}{2}$ 二阶相切（它们的二阶 Taylor 系数被标准化锁成同一个），而把这条共同的``种子''按 $n$ 次幂累积，极限就是高斯特征函数 $e^{-t^2/2}$。

\begin{figure}[h]
\centering
\begin{tikzpicture}[>=Stealth]
  \draw[->] (-2.6,0) -- (2.8,0) node[right] {$t$};
  \draw[->] (0,-1.0) -- (0,3.5) node[above] {$\vp(t)$};
  \node[anchor=north east] at (-0.05,-0.05) {\small$0$};
  \draw (-0.05,3.0)--(0.05,3.0) node[left,anchor=east,xshift=-2pt] {\small$1$};
  % 共同的二阶 Taylor 抛物线 1 - t^2/2
  \draw[thick,gray,densely dashdotted] plot coordinates
    {(-1.41,0.018) (-1.31,0.426) (-1.21,0.804) (-1.11,1.152) (-1.01,1.470)
     (-0.91,1.758) (-0.81,2.016) (-0.71,2.244) (-0.61,2.442) (-0.51,2.610)
     (-0.41,2.748) (-0.31,2.856) (-0.21,2.934) (-0.11,2.982) (-0.01,3.000)
     (0.09,2.988) (0.19,2.946) (0.29,2.874) (0.39,2.772) (0.49,2.640)
     (0.59,2.478) (0.69,2.286) (0.79,2.064) (0.89,1.812) (0.99,1.530)
     (1.09,1.218) (1.19,0.876) (1.29,0.504) (1.39,0.102)};
  % 某有限 n 的标准化特征函数
  \draw[thick,red!70!black,dashed] plot coordinates
    {(-2.20,0.078) (-2.00,0.198) (-1.80,0.391) (-1.60,0.657) (-1.40,0.989)
     (-1.20,1.367) (-1.00,1.765) (-0.80,2.152) (-0.60,2.497) (-0.40,2.767)
     (-0.20,2.940) (0.00,3.000) (0.20,2.940) (0.40,2.767) (0.60,2.497)
     (0.80,2.152) (1.00,1.765) (1.20,1.367) (1.40,0.989) (1.60,0.657)
     (1.80,0.391) (2.00,0.198) (2.20,0.078)};
  % 高斯特征函数 e^{-t^2/2}（极限）
  \draw[very thick,black] plot coordinates
    {(-2.20,0.267) (-2.00,0.406) (-1.80,0.594) (-1.60,0.834) (-1.40,1.126)
     (-1.20,1.460) (-1.00,1.820) (-0.80,2.178) (-0.60,2.506) (-0.40,2.769)
     (-0.20,2.941) (0.00,3.000) (0.20,2.941) (0.40,2.769) (0.60,2.506)
     (0.80,2.178) (1.00,1.820) (1.20,1.460) (1.40,1.126) (1.60,0.834)
     (1.80,0.594) (2.00,0.406) (2.20,0.267)};
  \node[gray,anchor=west] at (-2.55,-0.55) {\small$1-\tfrac{t^2}{2}$（共同二阶项）};
  \draw[gray,->] (-1.55,-0.45) -- (-1.32,0.34);
  \node[red!70!black,anchor=west] at (1.35,1.55) {\small$\vp_{Z_n}(t)$};
  \draw[red!70!black,->] (1.55,1.45) -- (1.45,0.92);
  \node[black,anchor=west] at (1.45,2.35) {\small$e^{-t^2/2}$};
  \draw[black,->] (1.7,2.25) -- (1.62,1.72);
\end{tikzpicture}
\caption{特征函数法的几何。标准化后，任意（有限二阶矩）分布的特征函数 $\vp_{Z_n}(t)$ 在 $t=0$ 处都与同一条抛物线 $1-\tfrac{t^2}{2}$ 二阶相切；把这个共同的二阶种子取 $n$ 次幂复利累积，极限就是高斯特征函数 $e^{-t^2/2}$（黑线）。纵轴已缩放，顶点对应 $\vp(0)=1$。}
\label{fig:clt-2}
\end{figure}

\rmkb[为什么不用矩母函数]{
矩母函数 $M_X(t)=\E{e^{tX}}$ 也能把独立和化为乘积，证明读起来更省力。但它的硬伤是\emph{不一定存在}——重尾分布（如 Cauchy、某些幂律收益）连均值都没有，矩母函数在 $t=0$ 之外可能发散。特征函数把 $e^{tX}$ 换成 $e^{itX}$，被积函数恒有界、期望永远存在，因此对一切分布都管用。这就是渐近统计宁用特征函数、不用矩母函数的原因。
}

\section{放松同分布：Lyapunov 与 Lindeberg}
\label{sec:clt-noniid}

iid 假设在计量里常常太强。回归里每个观测的解释变量取值不同，异方差使各项方差不等，面板与时间序列里各项的分布更是逐个相异。我们要的是：在``各项独立但\emph{不}同分布''时，标准化和何时仍趋于正态？关键直觉是——只要\emph{没有任何单独一项支配整个和}，正态性就保得住。把这句话定量化，就得到 Lindeberg 条件。

设 $X_1,\dots,X_n$ 独立（不必同分布），$\E{X_i}=\mu_i$，$\Var{X_i}=\sigma_i^2$，记总方差 $s_n^2=\sumi\sigma_i^2$。考察标准化和 $\dfrac{1}{s_n}\sumi(X_i-\mu_i)$。

\thm{Lindeberg 中心极限定理}{
若对每个 $\ve>0$ 都有\textbf{Lindeberg 条件}
\[
    \frac{1}{s_n^2}\sumi \E{(X_i-\mu_i)^2\,\indicator\set{\abs{X_i-\mu_i}>\ve s_n}}\ \longrightarrow\ 0
    \qquad(n\to\infty),
\]
则
\[
    \frac{1}{s_n}\sumi (X_i-\mu_i)\ \dto\ \cN(0,1).
\]
}

Lindeberg 条件读起来吓人，含义却朴素：方括号里的示性函数只挑出``偏离自身均值超过总标准差 $\ve s_n$''的那部分大涨落，条件要求所有这些大涨落对总方差的贡献\emph{一致地可忽略}。换句话说，没有哪一项（或哪几项）的尾巴大到能左右整个和——和是由许多势均力敌的小项堆出来的。这正是``众多小因素叠加''那句直觉的严格版本。

实践中 Lindeberg 条件不好直接验，于是常用一个更强、却好查的充分条件：

\cor{Lyapunov 中心极限定理}{
若存在某个 $\delta>0$ 使\textbf{Lyapunov 条件}
\[
    \frac{1}{s_n^{2+\delta}}\sumi \E{\,\abs{X_i-\mu_i}^{2+\delta}}\ \longrightarrow\ 0
    \qquad(n\to\infty),
\]
则 Lindeberg 条件成立，从而 $\dfrac{1}{s_n}\sumi(X_i-\mu_i)\dto\cN(0,1)$。
}

Lyapunov 条件只要求一个比方差略高的矩（取 $\delta=1$ 即三阶绝对矩、$\delta=2$ 即四阶矩）一致受控，便足以排除单项支配。它是计量实证里验证非同分布 CLT 最常用的入口——比如证明异方差下 OLS 估计量的渐近正态时，对误差项加一条``$2+\delta$ 阶矩有界''的矩条件，走的就是 Lyapunov 这条路。

\rmkb[iid 是非同分布版本的特例]{
当 $X_i$ iid 时，$\sigma_i^2\equiv\sigma^2$，$s_n^2=n\sigma^2$，$s_n=\sigma\rtn$，于是 $\frac{1}{s_n}\sum_i(X_i-\mu)=\frac{1}{\sigma\rtn}\sum_i(X_i-\mu)=\frac{\rtn(\bar X_n-\mu)}{\sigma}$——正好回到 Lindeberg–Lévy 的标准化形式。此时 Lyapunov 条件退化为``$\E{\abs{X_i-\mu}^{2+\delta}}<\infty$''（自动成立，因为 $\frac{n\cdot C}{(\sigma\rtn)^{2+\delta}}=\frac{C}{\sigma^{2+\delta}}n^{-\delta/2}\to0$）。所以 iid CLT 是这套更一般框架的一个干净特例。
}

\section{多元中心极限定理}
\label{sec:clt-multivariate}

计量里的估计量几乎总是\emph{向量}：一条回归有一整组系数，GMM 有一组矩条件。我们需要的是``随机向量的样本均值''的联合渐近分布——它决定了系数们\emph{一起}如何波动，也就决定了联合检验（Wald）的分布。所幸，多元版本可以由一元版本直接搭出来，靠的是一条把``向量收敛''归约到``所有标量方向收敛''的工具。

\thm{Cramér–Wold 投影法则}{
$\RR^k$ 中的随机向量列 $\set{\vc Z_n}$ 依分布收敛到 $\vc Z$，当且仅当对\emph{每个}固定方向 $\vc\lambda\in\RR^k$，标量投影都收敛：
\[
    \vc\lambda\T\vc Z_n\ \dto\ \vc\lambda\T\vc Z
    \qquad(\forall\,\vc\lambda\in\RR^k).
\]
}

这条法则的用处是``降维'':要判定一个向量列的依分布收敛，只需对它的每一条一维投影用我们已经会的一元 CLT。把它和 iid 一元 CLT 一拼，多元 CLT 立得。

\thm{多元中心极限定理}{
设 $\vc X_1,\vc X_2,\dots$ 为 $\RR^k$ 中的 iid 随机向量，均值 $\E{\vc X_i}=\vc\mu$，协方差矩阵 $\Var{\vc X_i}=\vc\Sigma$（各元素有限）。则
\[
    \rtn\,(\bar{\vc X}_n-\vc\mu)\ \dto\ \cN(\vc 0,\vc\Sigma),
\]
其中 $\bar{\vc X}_n=\frac1n\sumi\vc X_i$，$\cN(\vc 0,\vc\Sigma)$ 是均值 $\vc 0$、协方差 $\vc\Sigma$ 的多元正态。
}

\pf{
取任意固定方向 $\vc\lambda\in\RR^k$，把向量问题投影成标量。令 $Y_i=\vc\lambda\T\vc X_i$，则 $Y_1,Y_2,\dots$ 是 iid 标量随机变量，其均值与方差为
\[
    \E{Y_i}=\vc\lambda\T\vc\mu,
    \qquad
    \Var{Y_i}=\vc\lambda\T\vc\Sigma\,\vc\lambda .
\]
对这列标量用一元 Lindeberg–Lévy CLT：
\[
    \rtn\,(\bar Y_n-\vc\lambda\T\vc\mu)
    =\vc\lambda\T\,\rtn(\bar{\vc X}_n-\vc\mu)
    \ \dto\ \cN\!\pr{0,\ \vc\lambda\T\vc\Sigma\,\vc\lambda}.
\]
而 $\cN(0,\vc\lambda\T\vc\Sigma\vc\lambda)$ 恰是 $\vc\lambda\T\vc W$ 的分布，其中 $\vc W\dist\cN(\vc 0,\vc\Sigma)$（因为多元正态的线性投影 $\vc\lambda\T\vc W\dist\cN(0,\vc\lambda\T\vc\Sigma\vc\lambda)$，见『多维分布、独立性与变量变换』一章）。于是对每个 $\vc\lambda$ 都有 $\vc\lambda\T\rtn(\bar{\vc X}_n-\vc\mu)\dto\vc\lambda\T\vc W$。由 Cramér–Wold 法则，$\rtn(\bar{\vc X}_n-\vc\mu)\dto\vc W\dist\cN(\vc 0,\vc\Sigma)$。
}

这里浮现的协方差矩阵 $\vc\Sigma$ 就是后续一切``渐近协方差''的源头：它的对角元给出各分量的渐近方差（标准误的平方），非对角元给出分量间的渐近相关——正是联合假设检验需要的全部信息。

\section{去处：渐近推断的总枢纽}
\label{sec:clt-uses}

把估计量的渐近正态接回计量推断，是本章工具的归宿。下面这些是它的主战场。

\state{从 CLT 到一个置信区间}{
对 iid 样本，CLT 给 $\rtn(\bar X_n-\mu)\dto\cN(0,\sigma^2)$。用样本方差 $\hat\sigma^2\pto\sigma^2$ 替换未知的 $\sigma^2$（这一步替换的合法性由 Slutsky 定理保证，见后文『连续映射定理、Slutsky 与 Delta 法』一章），得 $t$ 形式的枢轴量
\[
    \frac{\bar X_n-\mu}{\hat\sigma/\rtn}\ \dto\ \cN(0,1).
\]
于是 $\mu$ 的渐近 $95\%$ 置信区间是 $\bar X_n\pm 1.96\,\hat\sigma/\rtn$。\emph{区间宽度以 $1/\rtn$ 收缩}——样本翻四倍，区间才缩一半，这条``$\rtn$ 速率''贯穿整个渐近推断。
}

\begin{itemize}
\item \textbf{估计量的渐近正态。}计量里的主流估计量（OLS、MLE、GMM、更一般的 M-估计量）几乎都能写成``某个 iid（或独立）量的样本平均''的形式，或经由一阶条件归约到这种形式。对那个平均用 CLT，就得到统一的结论 $\rtn(\hat{\vc\theta}-\vc\theta_0)\dto\cN(\vc 0,\vc V)$。这正是后文『极值估计的统一框架』一章的主线，那里的``三明治''协方差矩阵 $\vc V$ 就来自把 CLT 用在样本一阶条件上、再配上隐函数式的线性化。
\item \textbf{$t$、Wald、LM 检验。}一旦有了 $\rtn(\hat{\vc\theta}-\vc\theta_0)\dto\cN(\vc 0,\vc V)$，标准的检验统计量就都有了极限分布：单参数的 $t$ 比趋于 $\cN(0,1)$；多约束的 Wald 统计量 $\rtn(\hat{\vc\theta}-\vc\theta_0)\T\inv{\hat{\vc V}}\rtn(\hat{\vc\theta}-\vc\theta_0)$ 趋于 $\chi^2$（正态向量的二次型，见线性代数部分的二次型与谱定理）。这些临界值之所以``通用''，根子就在 CLT 把五花八门的数据统一压成了正态。
\item \textbf{Delta 法。}估计量常是某个渐近正态量的非线性变换 $\vc g(\hat{\vc\theta})$（弹性、比值、长期乘子……）。Delta 法用 $\vc g$ 的一阶 Taylor 把变换线性化，再借 CLT 的正态性得 $\rtn(\vc g(\hat{\vc\theta})-\vc g(\vc\theta_0))\dto\cN(\vc 0,\ \Jac_{\vc g}\vc V\Jac_{\vc g}\T)$。这正是下一章（『连续映射定理、Slutsky 与 Delta 法』）的内容；它整套建立在本章的渐近正态之上。
\item \textbf{自助法（bootstrap）的理论依据。}重抽样为何能逼近估计量的抽样分布？核心论证仍是 CLT：自助样本均值围绕原样本均值的标准化量，与 $\rtn(\bar X_n-\mu)$ 收敛到\emph{同一个}正态极限，于是用重抽样分布去近似真实抽样分布在一阶上成立。CLT 是自助法一致性的底层保证。
\end{itemize}

一条定理，一整套渐近推断。这是本书反复呈现的主题在统计这一端的化身：一个数学构造（标准化和趋于正态），被同一套逻辑反复用在每一个估计量、每一种检验、每一个置信区间上。认清了 CLT 这把钥匙，渐近统计后续诸章——Slutsky、Delta 法、极值估计的统一框架——便都成了对它的展开与精修。
